\section{Actual cubic Type II tangent rigidity in the distinguished sectors}
\label{sec:cubic-typeii-rigidity}

We use the actual complete cubic ring constructed in
AC-G49--G54 and its selected local component.  In this section
$K=\mathbb Q(\omega)$, $\eta=\eta_K$, $\chi=\chi_{269}$,
$j\in\{1,2\}$, $\gamma=\psi\kappa_\ell^e$, and
$e\ge1$.  All old \texttt{GD} and local published hypotheses
are retained.  The source paths have the pinned root
\path{Development/AutomorphicArithmetic/AutomorphicVarieties/}
at commit \texttt{7cbf7a0c32fca1badf05931879a37c3f8a2b03b4}.
No assertion here concerns Type I points.

\begin{lemma}[AC-G62-cont: every actual dual point is continuous]
\label{AC-G62-cont}
Let $O$ be the integers of a finite extension of
$\mathbb Q_\ell$, and let $S$ be a complete Noetherian
local $O$-algebra with the source's original residue field.
Let
\[
 F:S\longrightarrow_O\overline{\mathbb Q}_\ell
                         [\varepsilon]/(\varepsilon^2),
       \qquad F=f+\varepsilon d,                         \tag{G62.1}
\]
be any algebra map with $f$ an actual local point.
There is one finite extension $L/\operatorname{Frac}O$
containing both components of $F(S)$; its image is a finite
$O$-module.  The two components are continuous for the
maximal-ideal topology on $S$ and the valuation topology on
$L$.  Thus every such map is an actual continuous finite
coefficient dual point, without a continuity assumption on
$d$ and without \texttt{M2}.
\end{lemma}
\begin{proof}
The automatic-locality and finite-image proof of
AC-G58-algebra applies to the complete local ring $S$.
Explicitly, a reversed monic equation for $f(x)^{-1}$
would contradict the invertibility of $1+\mathfrak m_S$
if $x\in\mathfrak m_S$ and $|f(x)|_\ell\ge1$.
Finite generators of $\mathfrak m_S$ then give a
contraction bound for its powers.  Residue-field
approximation by polynomials in those generators shows
that $f(S)$ lies in the integers of a single finite
extension, and hence is a finite $O$-module.  This is
the general local-ring part of the proof of (G58.1),
not a use of degree-zero global finiteness.

Put $I=\ker f$.  Since $S$ is Noetherian, $I/I^2$ is
finite over $S/I=f(S)$.  The exact sequence
\[
 0\longrightarrow I/I^2\longrightarrow S/I^2
                         \longrightarrow S/I\longrightarrow0
                                                               \tag{G62.2}
\]
therefore makes $S/I^2$ a finite $O$-module.
For $x,y\in I$, multiplicativity gives
$F(xy)=\varepsilon^2d(x)d(y)=0$.  Thus $F$ factors
through this finite module.  Images of finitely many
$O$-generators of $S/I^2$ put both components in a
single finite extension $L$, and give a bounded finite
$O$-module containing $F(S)$.

To verify continuity directly, choose generators
$u_1,\ldots,u_t$ of $\mathfrak m_S$, and put
\[
 \rho=\max(|\varpi|_\ell,|f(u_1)|_\ell,\ldots,
                         |f(u_t)|_\ell)<1.
\]
Choose a constant $C$ bounding both components of
$F(S)$.  The product rule on a product of $n$ of the
$u_i$ gives a first component of norm at most
$\rho^n$ and a second component of norm at most
$C\rho^{n-1}$.  Multiplication by an arbitrary
coefficient from $S$ adds at most a fixed further
factor to this bound.  The ultrametric inequality
handles arbitrary sums without a factor depending
on their number.  Consequently, after increasing
$C$ once,
\[
 |f(x)|_\ell\le C\rho^n,\qquad
 |d(x)|_\ell\le C\rho^{n-1}
                  \quad(x\in\mathfrak m_S^n).          \tag{G62.3}
\]
Both functions are continuous.  This proves the
lemma for every lift (G62.1), not only for derivations
first specified on a dense subalgebra.
\end{proof}

\begin{lemma}[AC-G62-block: actual Type II middle and determinant]
\label{AC-G62-block}
Let $\ell\ne2,3,269$ and let $r_f$ be an actual
local-point specialization of the cubic universal ring,
with Type II endpoint block $B$ on $G_K$ and middle
character $m$.  Use the exact ordinary inertia
polynomial supplied by AC-G49-inert or AC-G57 at
$\ell$, and retain the original residual-separation
condition \texttt{GD.gen}.  Then
\[
                   m=1,\qquad\det B=\chi^{2j}.          \tag{G62.4}
\]
No unramifiedness of $\psi$ at $269$ is required.
The endpoint untwist is an actual old-model point
as in AC-K29-point and AC-G48-cubic, with all the
old point-transport hypotheses retained.
\end{lemma}
\begin{proof}
By AC-G62-cont applied to the constant dual point,
the representation has values in one finite extension.
Its values are integral and its specified reduction
on $G_K$ is
\[
   \operatorname{diag}(\overline\chi^{\,j}\overline\gamma,
                  1,\overline\chi^{\,j}\overline\gamma^{-1}).
                                                               \tag{G62.5}
\]
For a finite place $v\nmid\ell$ of $K$, the actual
universal killing condition makes $r_f|_{I_v}$
factor through the image of (G62.5) on $I_v$.
That image is a finite subgroup of a diagonal torus
over the finite residue field, hence has order prime
to $\ell$.  The character $m|_{I_v}$ is therefore
of prime-to-$\ell$ finite order and has reduction
one.  Reduction is injective on such roots of
unity, so $m|_{I_v}=1$.

The same argument applies to the determinant.
If $d$ is the order of the image of an inertia
element in (G62.5), its actual image has $d$th
power one, and $\chi^{2jd}=1$ as well: this last
identity follows by reduction and $\ell\ne3$.
Thus the ratio $\det(r_f)\chi^{-2j}$ has finite
prime-to-$\ell$ inertia image and reduction one,
so is trivial on $I_v$.  This argument includes
arbitrary simultaneous ramification of $\psi$
and $\chi$ at $269$.

At a place above $\ell$, $\chi$ is unramified,
and the actual inertia characteristic polynomial
has roots $\gamma(h),1,\gamma(h)^{-1}$.
Whenever $\overline\gamma(h)\ne1$, the middle
character, whose reduction is one, must equal
one.  Choose $s$ in inertia with
$\overline\gamma(s)\ne1$, as supplied by
\texttt{GD.gen}.  If $\overline\gamma(h)=1$,
both $s$ and $hs$ have nontrivial residual
$\gamma$-value, so
$m(h)=m(hs)m(s)^{-1}=1$ also.
The determinant polynomial gives
$\det r_f|_{I_\ell}=1$ directly.

Both $m$ and $\det(r_f)\chi^{-2j}$ are now
everywhere-unramified characters of $G_K$.
The imaginary quadratic field $K$ has class
number one, so class field theory makes both
characters trivial.  This proves (G62.4).

Apply AC-G45 to extend the determinant-one
endpoint untwist to an odd rational representation
$V_f$ with determinant $\eta_K$.
The exact modified-kernel-to-old-kernel argument
of AC-K29-kill applies, including when $\psi$
is ramified at $269$.  At the coefficient prime,
the actual local component supplies its ordered
ordinary flag and crystallinity on the source's
$H_L$; for the new split component this is the
fixed type of AC-G54 and AC-G57, and for the
transported component it is the old local
condition of AC-G49-inert.
In the fixed-type description, restriction to
$H_L$ kills the prescribed finite inertial type;
zero monodromy then means crystallinity there.
AC-G48-cubic gives precisely (GP.11).
AC-K29-point therefore supplies the actual old
universal-ring point, with its original residue
field and normalization.  These are pointwise
transports; no transport of the mixed dual
deformations into the old ring has been used.
\end{proof}

\begin{theorem}[AC-G62: actual Type II rigidity]
\label{AC-G62}
Retain the actual cubic model of AC-G54-model,
the entire old \texttt{GD} and \texttt{CohInp},
the old pointwise input
\path{AutomorphyUnramifiedCoefficientsAlt} with
its coefficient restriction $e_O=1$, and the
literal old \texttt{RigiditySecondCase}.
Suppose $\ell\ne2,3,269$ and either
\[
 \ell\text{ is inert in }K,\qquad\text{or}\qquad
 \ell=\mathfrak p\overline{\mathfrak p}
   \text{ and }\chi(\operatorname{Frob}_{\mathfrak p})\ne1.
                                                               \tag{G62.6}
\]
For every actual Type II local point
$f:S\to\overline{\mathbb Q}_\ell$ of the cubic ring,
and every algebra lift $F$ in (G62.1) with first
component $f$, there is a matrix $X$ such that
\[
       \operatorname{tanCoc}(F,r^{\rm univ})(g)
             =X-\operatorname{ad}(r_f)(g)X
                         \quad\text{for every }g.       \tag{G62.7}
\]
This is the literal Type II part of
$\operatorname{Rigid}(r^{\rm univ},f)$.
In fact, for every finite coefficient field $L$
containing the matrices of $r_f$, one has
\[
             H_f^1(\mathbb Q,\operatorname{ad}_L r_f)=0.
                                                               \tag{G62.8}
\]
The quantifiers include every coefficient extension,
marked element and ring/model choice allowed by
the stated construction.  No \texttt{M2}, Type I
rigidity, or integral Selmer-length estimate is
deduced.
\end{theorem}

\begin{proof}
\proofstep{1. The actual tangent lies in the stated local subspace.}
AC-G62-cont puts every $F$ in one finite
coefficient dual-number algebra and proves its
continuity.  Writing $F=f+\varepsilon d$, the
actual tangent is
\[
 z_F(g)=d(r^{\rm univ}(g).\mathrm{mat})
                   f(r^{\rm univ}(g).\mathrm{mat})^{-1}.
                                                               \tag{G62.9}
\]
The product rule and the differential
$\theta_{\rm Lie}(Y)=-JY^tJ$ prove the cocycle
identity, including the nonidentity semidirect
component.  Matrix inversion on the invertible
locus is continuous, so this is a continuous
finite-dimensional cocycle.

The actual local-condition assertion is already
AC-G53-BK.  We recall its precise use here.
Away from $\ell$, the literal modified killing
condition makes (G62.9) zero on inertia intersected
with the actual cubic residual kernel.  Its
restriction factors through a finite quotient;
averaging makes it an inertia coboundary over
the coefficient field.  At a transported local
component use the full old dual-point theorem
after the constant conjugation of AC-G49-inert.
At the new split component, every finite local
dual point killing $P_*$ is potentially crystalline
by AC-G54(3).  Over a finite extension on which
the type is crystalline, its underlying
self-extension becomes split after tensoring
with $B_{\rm cris}$: the period dimension is
twice that of the base representation.
Restriction--corestriction injects rational
$H^1$ with these period-ring coefficients,
because the finite extension degree is invertible.
Thus it is a period-ring coboundary already over
the rational local group.  This is exactly
$z_F\in H_f^1$ at the cocycle level.
Only this one-way assertion about actual dual
points is needed; no equality with an entire
abstract tangent space is assumed.

\proofstep{2. Transport the point and its corner, not its mixed deformation.}
AC-G62-block gives the middle and determinant
identities and the actual old-model point
$\widetilde r_f$ after endpoint untwisting.
The retained old automorphy input and
AC-K29-typeII supply its rational representation
$V_f$, its actual newform $f_V$ of weight
$2e+1$, its original level conditions and
$\operatorname{IsModularDense}$.
Its actual ordinary flag is the one used in
AC-K29-nonCM.  That theorem excludes every
quadratic self-twist, so $f_V$ is non-CM.
AC-K29-split supplies its characteristic-zero
derived-image conclusion.  None of these
assertions is inferred from a trace-only
comparison with an abstract representation.

Use the three equivariant coordinate projections
$P_A,P_C,P_R$ in (GT.3)--(GT.4).
Endpoint untwisting multiplies the corner
matrices by a scalar.  Scalars cancel in
conjugation, and the inverse-transpose action
on the other coset is unchanged.  Hence the
corner adjoint for $r_f$ is literally the
corner adjoint for $\widetilde r_f$.

The exact retained source input is
\path{ModularityArgument/ThresholdArguments/DenseThresholdCases/RigEnd3.lean},
\texttt{RigiditySecondCase}.  It quantifies over
the old datum, an actual old ring and local
point, its Type II predicate, every continuous
$W$ satisfying \texttt{IsExt}, determinant
$\eta_K$, and its actual
\texttt{IsModularDense}; it concludes
\texttt{AdZero2}.  All these hypotheses are
supplied by the point transport just described.
The corner application of \texttt{AdZero2}
in the complete proof of
\path{ModularityArgument/ThresholdArguments/DenseThresholdCases/TypeIIOfBlochKatoLocalPoint.lean}
uses the explicit symmetric conjugate-dual pair
$J_2$, $V_f(c)=\operatorname{diag}(-1,1)$,
and the equivariant maps \texttt{pA},
\texttt{sA}.  These data are unchanged here.
It therefore gives
\[
                      H_f^1(\mathbb Q,\mathcal A_f)=0.
                                                               \tag{G62.10}
\]
This use is solely the corner step of that proof;
its old \texttt{kUse'} mixed-character step is
not invoked.

The middle adjoint is the one-dimensional
finite character $\eta_K$: the inverse-transpose
differential changes the middle entry's sign
on the nontrivial coset.  Thus
AC-G53-Artin with weight zero gives
\[
                      H_f^1(\mathbb Q,\mathcal C_f)=0.
                                                               \tag{G62.11}
\]

\proofstep{3. The newly proved mixed vanishing.}
The precise remaining coefficient is, by
AC-G45-cohomology,
\[
 \mathcal R_f\simeq V_f\otimes
             \operatorname{Ind}_{G_K}^{G_{\mathbb Q}}\chi^j.
                                                               \tag{G62.12}
\]
At an inert prime, all its actual ordinary,
potentially crystalline, modular and
derived-image hypotheses were checked above,
so AC-K29-inert gives
$H_f^1(K,V_f\chi^j)=0$.
At a split prime in (G62.6), use
AC-K29-primitive and AC-K29-ramified; the
source's \texttt{GD.ine} is precisely the
quadratic conductor-component condition needed
for the latter, including the trivial case.
These theorems apply to the same original
$V_f$, rather than only to its good ordinary
quadratic twist.  Their finite cyclotomic
specialization cancels that auxiliary twist.
Induction preserves the specified Bloch--Kato
local kernels, as proved in AC-G45-cohomology.
Consequently
\[
                      H_f^1(\mathbb Q,\mathcal R_f)=0.
                                                               \tag{G62.13}
\]

Work over a finite field containing the point,
the chosen newform coefficients and the actual
cocycle.  Such a field exists for every lift
by AC-G62-cont.  The assertions of rational
cohomology commute with a finite coefficient
extension, so their vanishing descends to any
smaller finite field of definition as well.
The three projections are equivariant and
sum to the identity; they preserve inertia
and period-ring coboundaries.  Equations
(G62.10)--(G62.13) therefore prove (G62.8).
Applied to the actual cocycle of Step~1 they
produce three cobounding matrices.  Their
sum, with the source convention
$\operatorname{cob}(r,X)=X-\operatorname{ad}(r)X$,
gives (G62.7).  This works for every actual
lift $F$, so it is precisely the required
Type II \texttt{Rigid} predicate and kills
the actual image $\mathcal T_f$ in (GT.5).
\end{proof}

\paragraph{Exact scope and geometric dictionary.}
The new arithmetic use is the vanishing of the
actual mixed Type II image in the sectors (G62.6).
The continuity lemma removes an unnecessary
extra premise on algebraic dual points; the
local Bloch--Kato conclusion itself is the
already proved AC-G53-BK.  The old corner
rigidity and pointwise modularity remain the
literal inputs named in the theorem, including
the coefficient-ring and \texttt{GD} conditions.
Type I points, split irregular CM primes,
the exceptional primes $2,3,269$, global
\texttt{M2}, and the numerical integral index
bound remain separate.  The finite Artin
characters are finite local systems and the
block projections are morphisms of their
$\ell$-adic realizations in the already
audited SGA~1, Expos\'e~V, \S\S4 and~7
dictionary.  The vanishing arguments come
from the precisely stated arithmetic inputs,
not from that categorical interpretation.

\begin{corollary}[AC-G62-all: actual Type II rigidity at every good coefficient prime]
\label{AC-G62-all}
Retain the actual cubic model of AC-G54-model, the entire
old \texttt{GD} and \texttt{CohInp}, the literal old
\texttt{AutomorphyUnramifiedCoefficientsAlt} input including
its coefficient restriction $e_O=1$, and the literal old
\texttt{RigiditySecondCase}.  For every prime
$\ell\notin\{2,3,269\}$, every actual Type II local point
$f$ of that cubic model, and every algebra lift $F$ with
first component $f$, the identity (G62.7) holds.
For every finite coefficient field $L$ containing the
matrices of $r_f$, one also has (G62.8).
There is no residual CM-distinction condition at a split
coefficient prime.
\end{corollary}
\begin{proof}
AC-G62-cont and AC-G62-block apply to every prime in
the stated range.  They give a continuous finite-coefficient
dual point, the actual middle and determinant identities,
and the actual old-model point after endpoint untwisting.
The local Bloch--Kato argument and the complete corner and
middle arguments in Steps~1--2 of AC-G62 therefore apply
unchanged, under exactly the inputs retained here.

If $\ell$ is inert in $K$, Step~3 is already
AC-K29-inert.  If $\ell$ is split, apply
AC-K29-irregular: its actual local, non-CM and
characteristic-zero derived-image hypotheses are supplied
by the same point transport, ordinary flag,
AC-K29-nonCM and AC-K29-split used in Step~2.
Its quadratic cyclotomic specialization, when necessary,
returns the original $V_f$.  The theorem gives
$H_f^1(K,V_f\chi^j)=0$ for both $j=1,2$, without a
condition on $\chi(\operatorname{Frob}_{\mathfrak p})$.
Thus (G62.13) holds in every split case as well.

The same three equivariant coordinate projections now
give (G62.8) over any finite coefficient field by finite
coefficient descent.  Applied to each actual tangent,
their three cobounding matrices sum to the matrix $X$
in (G62.7).  This proves the literal Type II
\texttt{Rigid} predicate for every lift, with all
original marked-element, datum and model quantifiers.
\end{proof}

\begin{remark}[Scope after AC-G62-all]
The split irregular exclusion in the preceding scope
paragraph is removed by this corollary.  The exceptional
primes $2,3,269$, the original coefficient-ring and
\texttt{GD} restrictions, Type I rigidity and Hecke
support, integral \texttt{M2}, and the quantitative
integral index remain separate.  The new conclusion is
rational Type II rigidity; the retained old modularity
and corner-rigidity inputs have not been generalized.
\end{remark}
